\section*{第 \ref{ch:b2:wave-interfaces} 章——界面处的反射与透射}

\begin{solution}{exo:b2:wave-interfaces:1}
$r = (Z_2 - Z_1)/(Z_2 + Z_1)$，$R = r^2$：空气–水 $0.9995$，$99.9\%$；
水–钢 $0.93$，$86\%$；组织–骨 $0.63$，$40\%$；组织–空气
$-0.9995$，$99.9\%$。你在空气中的声音被表面反射；
引擎使船体振动，船体是与水接触的钢，并
将其 $14\%$ 的振动直接透入水中。
\end{solution}

\begin{solution}{exo:b2:wave-interfaces:2}
$R = ((n - 1)/(n + 1))^2$：$2.0\%$，$4.3\%$，$17\%$；玻璃–水 $0.44\%$。
$0.957^{10} = 0.64$：三分之一的光损失了，其中大部分成了鬼影；
镀膜后：$0.995^{10} = 0.95$。
\end{solution}

\begin{solution}{exo:b2:wave-interfaces:3}
$\tan\theta_B = n_2/n_1$：\qty{56.7}{\degree}，\qty{53.1}{\degree}，\qty{48.8}{\degree}。
反射光偏振方向垂直于入射面；
透射光只是部分偏振（s 分量
只是部分被反射）。在与法线成 \qty{57}{\degree} 角处观察窗户，
并使起偏器的轴位于入射面内（
对从侧面看的竖直窗户而言为水平方向）。
\end{solution}

\begin{solution}{exo:b2:wave-interfaces:4}
$\arcsin(1/n)$：\qty{41.1}{\degree}，\qty{48.8}{\degree}，\qty{24.4}{\degree}；$45^\circ > 41^\circ$：
棱镜发生\omterm{thm:b2:wave-interfaces:snell}{全反射}。$\ell = \lambda/2\pi\sqrt{n^2\sin^2\theta - 1}$：
\qty{45}{\degree}时为 \qty{270}{nm}，\qty{60}{\degree}时为 \qty{115}{nm}。金刚石较小的临界角
使进入其中的光被困住，直到从顶面射出——
即钻石的璀璨光彩。
\end{solution}

\begin{solution}{exo:b2:wave-interfaces:5}
(a) $\sqrt{1.52} = 1.23$；$e = 550/4 \times 1.38 = \qty{100}{nm}$。(b) $r_1 = -0.160$，$r_2 =
-0.048$，$\varphi = \pi$：$r = (r_1 - r_2)/(1 - r_1r_2) = -0.112$，$R = 1.3\%$（相较于
$4.3\%$）。(c) \qty{400}{nm}：$\varphi = 1.375\pi$，$R \approx 2.2\%$；\qty{700}{nm}：$\varphi =
0.79\pi$，$R \approx 1.6\%$：蓝光和红光的反射比绿光更强——
即镀膜透镜呈现的紫色光泽。(d) 多层膜可使 $R$
在整个可见光范围内趋于平坦。
\end{solution}

\begin{solution}{exo:b2:wave-interfaces:6}
(a) $T = 4 \times 30 \times 1.6/31.6^2 = 19\%$；其余部分在层内往复反射，脉冲
持续振铃。(b) $Z = \sqrt{Z_1Z_2} = \qty{6.9e6}{kg.m^{-2}.s^{-1}}$，$e = v/4f = \qty{0.15}{mm}$。
(c) $T_1 = 4 \times 30 \times 10^6 \times 410/(30 \times 10^6)^2 = 5.5 \times 10^{-5}$，$T_2 = 4 \times 1.6 \times
10^6 \times 410/(1.6 \times 10^6)^2 = 1.0 \times 10^{-3}$：$5.6 \times 10^{-8}$。(d) $4 \times 1.5 \times 1.6/
3.1^2 = 0.999$。
\end{solution}

\begin{solution}{exo:b2:wave-interfaces:7}
(a) $\omega_p/\omega = 3$：$\underline n = -\iu\sqrt8 = -2.83\iu$；$r = (1 + 2.83\iu)/(1 - 2.83
\iu)$，$|r| = 1$，相位 $2\arctan2.83 = \qty{141}{\degree}$。(b) $c/\sqrt{\omega_p^2 - \omega^2}
= \qty{5.6}{m}$。(c) $E \propto \cos(kx + \varphi/2)$：第一个波节位于 $kx = -(\pi + \varphi)/2$，
$x \approx \qty{-45}{m}$，位于该层下方。(d) 折射率随高度下降，因此
笛卡儿使光线偏离竖直方向；光线在
$n(z) = \sin\theta_0$ 处折返，即 $\omega_p^2 = \omega^2\cos^2\theta_0$，位于 $\omega = \omega_p$ 所对应高度之下；
若频率的 $\omega\cos\theta_0$ 超过该层的最大值 $\omega_p$，
则该波逸出。
\end{solution}

\begin{solution}{exo:b2:wave-interfaces:8}
$\sqrt{8\varepsilon_0\omega/\gamma}$：$2.7 \times 10^{-4}$（$R = 0.9997$），$1.5 \times 10^{-2}$（$0.985$），
$0.06$（$0.94$）——但在 \qty{5e14}{Hz} 时 $\omega\tau = 80$：铜是\omterm{def:b2:waves-in-media:plasma}{等离子体}，
不是欧姆导体，约反射 $60\%$，且有颜色。\qty{15}{W}
的功率被 CO$_2$ 光束吸收：在
小镜子上每分钟升温几百开尔文，因此需用水冷却。
\end{solution}

\begin{solution}{exo:b2:wave-interfaces:9}
$r = 0.2$，$-1/3$，$-1$，$+1$，$0$；\omterm{prop:b2:wave-interfaces:coax}{驻波比} $1.5$，$2$，$\infty$，$\infty$，$1$；
被反射的比例分别为 $4\%$，$11\%$，$100\%$，$100\%$，$0$。\qty{96}{W} 到达天线，\qty{4}{W}
返回发射机（并会部分再次被反射）。\qty{1.2}{\micro s}：
\qty{120}{m}；正相回波：断开的电路——电缆被切断。
\end{solution}

\begin{solution}{exo:b2:wave-interfaces:10}
(a) $\vect E = E\vect e_z$；$\vect k\wedge\vect e_z = (k_y, -k_x, 0)$，因此切向
分量在入射波与透射波中为 $B_y = -(n/c)\cos\theta\,E$，在
反射波中为 $+(n_1/c)\cos\theta_1E_r$。(b) $1 + r = t$ 且 $n_1\cos\theta_1
(1 - r) = n_2\cos\theta_2t$：即为题中所给的 $r_s$。(c) $\theta = 0$：$(n_1 - n_2)/(n_1 +
n_2)$；$\theta_1 \to 90^\circ$：$r_s \to -1$。(d) $\cos\theta_2 = 0.882$：$r_s = -0.30$，$R_s =
9.2\%$；$R_p = 0.85\%$；与垂直入射时的 $4.3\%$ 相比。
\end{solution}

\begin{solution}{exo:b2:wave-interfaces:11}
(a) \qty{270}{nm}。(b) $d = 0.35\ell = \qty{95}{nm}$；$2.3\ell = \qty{620}{nm}$。(c)
脊与玻璃接触，使光射入皮肤：
探测器上脊对应暗纹，谷对应亮纹。(d) \omterm{def:b2:dispersion-wave-packets:complex}{倏逝波}是
禁区中按指数衰减的波函数；衰减长度内的
第二种介质将其收集——光学隧道效应。
\end{solution}

\begin{solution}{exo:b2:wave-interfaces:12}
(a) 两镜面处 $\vect E$ 的波节：$\sin kL = 0$，$L = m\lambda/2$，$f_m = mc/2L$。
(b) $1$，$2$，\qty{3}{GHz}：\qty{2.45}{GHz} 处无任何模式——实际微波炉的
三维模式要密集得多。(c) 每个镜面上 $j_s = H_0 = E_0/\mu_0c = \qty{53}{A/m}$，
且压强 $\mu_0j_s^2/2 = \qty{1.8}{mPa}$ 将镜面推开。
$\delta = \qty{1.3}{\micro m}$，$R_s = \qty{12.5}{m\ohm}$；$\tfrac12R_sH_0^2 =
\qty{18}{W/m^2}$，每个镜面；平均储能 $\tfrac14\varepsilon_0E_0^2 \times L = \qty{1.3e-7}{J/m^2}$；
$Q = 2\pi f \times 1.3 \times 10^{-7}/35 \approx 6 \times 10^4$。
\end{solution}

\begin{solution}{pb:b2:wave-interfaces:1}
\textbf{1.} $R = 4.3\%$；$0.957^{16} = 0.50$；一半的光成为杂散光。

\textbf{2.} $n = 1.23$，$e = \qty{112}{nm}$；没有耐用的固体具有如此低的
折射率（MgF$_2$，$1.38$，为最低值）：残余反射率 $1.3\%$。

\textbf{3.} $0.987^{16} = 0.81$；$0.997^{16} = 0.95$。杂散反射使
图像蒙上雾霾：遮蔽性眩光，对比度低。

\textbf{4.} $\varphi = \pi \times 550/440 = 1.25\pi$：$R \approx 0.0256 + 0.0025 + 2 \times 0.008 \times
\cos1.25\pi = 1.7\%$——比 \qty{550}{nm} 处更大，与红端情形类似：一个
品红色反射。

\textbf{5.} 每个界面的\omterm{thm:b2:wave-interfaces:normal}{反射系数}模为 $|r| = 0.97/3.73 = 0.26$；四分之一波长
间距使相继的反射同相，因此振幅相加，
而不是相互抵消；叠层的反射率为 $1 - 4(n_L/n_H)^{2N}
\approx 1 - 10^{-9}$，对 $N = 20$。

\textbf{6.} 同相条件仅在设计
波长附近成立（约 $\pm15\%$ 的频带）；其余波段透过：绿色镜子
透射时呈品红色。

\textbf{7.} 能量转入透射波：干涉使能量重新分配，
$R + T = 1$。

\textbf{8.} $19\%$；取 $Z_m = \sqrt{Z_1Z_2} = \qty{6.9e6}{kg.m^{-2}.s^{-1}}$，在
\qty{5}{MHz} 处为 $100\%$。

\textbf{9.} $r_1 = 0.63$，$r_2 = 0.62$；在 \qty{3}{MHz} 处往返相位为
$0.6\pi$，在 \qty{7}{MHz} 处为 $1.4\pi$：$R \approx 0.39 + 0.39 - 0.24 = 0.54$，$T \approx 0.46$
两处相同：匹配仅在 \qty{5}{MHz} 附近的频带内成立，而
短脉冲需要宽频带——这是一种折中。

\textbf{10.} $t_1 = 2 \times 410/30 \times 10^6 = 2.7 \times 10^{-5}$，$t_2 = 2.0$：$t = 5.5 \times 10^{-5}$，
$T = t^2Z_1/Z_3 = 5.6 \times 10^{-8}$：\qty{-72}{dB}。

\textbf{11.} $t = 0.095 \times 1.03 = 0.098$，$T = 0.098^2 \times 30/1.6 = 0.18$：凝胶
仅仅排除了空气。

\textbf{12.} $0.9^N = 0.01$：$N = 44$ 次往返——余振很长；
背衬吸收反向波，将其消除。

\textbf{13.} 骨骼处 $R = 0.40$；$2 \times 5 \times 0.5 \times 5 = \qty{25}{dB}$ 的衰减，即
$0.40 \times 3.2 \times 10^{-3} = 1.3 \times 10^{-3}$（\qty{-29}{dB}）；骨骼
吸收未被反射的部分：骨骼之后没有任何回波返回。

\textbf{14.} $R = 1$，$r = -1$。

\textbf{15.} 在 $x = 0$ 处短路：$v = -2\iu v_i\sin kx$，$i = (2v_i/Z_0)\cos kx$；
在 $x = -d$ 处：$v/i = \iu Z_0\tan kd$；当 $d = \lambda/4$ 时为无穷大。

\textbf{16.} $R_s \parallel \infty = R_s$：$r = (R_s - Z_0)/(R_s + Z_0) = 0$，当
\qty{377}{\ohm} 时；$d = \qty{7.5}{mm}$。

\textbf{17.} 薄片感受到全部入射场：$|v|^2/2R_s =
E_0^2/2Z_0$，即入射强度——全部被吸收。

\textbf{18.} 在 \qty{8}{GHz} 处：$kd = 72^\circ$，$Z_{\text{入}} = \iu\qty{1160}{\ohm}$，$R_s \parallel
Z_{\text{入}} = (340 + 111\iu)\,\Omega$，$R = |r|^2 = 2.6\%$（\qty{-16}{dB}）；在 \qty{12}{GHz} 处：
$kd = 108^\circ$，模相同：\qty{-16}{dB}。在约 $\pm20\%$ 的频带内有效。

\textbf{19.} 斜入射使有效的四分之一波长缩短
（$d\cos\theta$），并且波阻抗随角度改变：
匹配关系被破坏。

\textbf{20.} 单层屏只在一个频率和一个角度下起作用；渐变的
有耗层与成形表面则能在宽频带内吸收并偏转波。

\textbf{21.} $\underline n = -\iu\sqrt{(\omega_p/\omega)^2 - 1} = -3.6\iu$；$|r| = 1$，相位
$2\arctan3.6 = \qty{149}{\degree}$；穿透深度 $c/3.6\omega = \qty{22}{nm}$。波
从不进入金属，因而不会分裂为两个相当的反射：
\omterm{prop:b2:wave-interfaces:optics}{四分之一波长层}没有可与模为一的反射相抵消的
对象。

\textbf{22.} \qty{30}{W}；$30/(0.01 \times 235) = \qty{13}{K/s}$。

\textbf{23.} 每隔 $\lambda/2 = \qty{250}{nm}$ 为一个波节；$E \propto \cos(kx + \varphi/2)$，其中 $\varphi =
149^\circ$：第一个波节位于 $x = -(90^\circ + 74.5^\circ)\lambda/360^\circ = \qty{-230}{nm}$；
该处的晶粒处于暗处。

\textbf{24.} $I = \qty{1e7}{W/m^2}$，$E_0 = \sqrt{2I/\varepsilon_0c} = \qty{87}{kV/m}$，$j_s =
2E_0/\mu_0c = \qty{460}{A/m}$。

\textbf{25.} 金属：表\omterm{prop:b2:maxwell-equations:boundary}{面电流}位于\omterm{def:b2:dispersion-wave-packets:complex}{趋肤深度}之内，$r \approx -1$，
受限于这些电流的焦耳损耗。\omterm{def:b2:waves-in-media:plasma}{等离子体}：倏逝穿透，
$|r| = 1$，相位随频率而变，受限于碰撞。
介质叠层：许多弱反射同相叠加，$r \to 1$ 出现在
一个频带内，受限于层数和带宽。
\end{solution}
